\subsection{Brahman, Atman, and Maya}

The plain picture now receives its technical vocabulary. The Hindu traditions, from the Upanishads through Vedanta and the Bhagavad Gita, name the ground of all being, the innermost self, and the power by which the one appears as the many. They do not agree on which of these is primary. That disagreement is one of the tradition's most useful features for comparison.

\subsection{Brahman, the one ground}

\textbf{Brahman} is the one ground of all being, the single reality behind everything. It is described in two registers at once.

\begin{itemize}
 \item Brahman \textbf{without qualities} (\textbf{Nirguna Brahman}). Beyond every attribute, name, form, and thought. The pure undivided reality, reached only by negation. The classic method is "not this, not this" (\textbf{neti neti}). For the non-dualist school this is the true and only absolute.
 \item Brahman \textbf{with qualities} (\textbf{Saguna Brahman}). The absolute as it becomes knowable, the object of worship and meditation, the creator and Lord. This register is often identified with the personal God (\textbf{Ishvara}).
\end{itemize}

The classic formula for Brahman's nature is being, consciousness, and bliss (\textbf{sat-chit-ananda}).

\subsection{Atman and the great sayings}

\textbf{Atman} is the \textbf{self}, the innermost consciousness, the witness at the center of a person. The decisive Upanishadic claim is that Atman and Brahman are not two. The self at the center is identical with the ground of all being. This is condensed in the great sayings (the \textbf{Mahavakyas}).

\begin{longtable}{@{} L{0.30\textwidth}L{0.30\textwidth}L{0.30\textwidth}@{}}
\caption{The four Mahavakyas, the great sayings of identity.}\\
\toprule
\textbf{Saying} & \textbf{Meaning} & \textbf{Source} \\
\midrule
Tat tvam asi & that thou art & Chandogya Upanishad 6.8.7 \\
Aham brahmasmi & I am Brahman & Brihadaranyaka Upanishad 1.4.10 \\
Prajnanam brahma & consciousness is Brahman & Aitareya Upanishad 3.3 \\
Ayam atma brahma & this self is Brahman & Mandukya Upanishad 2 \\
\bottomrule
\end{longtable}

\subsection{Maya, the power of appearance}

\textbf{Maya} is the power by which the one Brahman appears as the many. Its meaning shifts by school. In the non-dualist reading, maya is the cosmic illusion. The world is not flatly unreal, but it is not finally real either. It is dependent appearance (\textbf{mithya}). Only Brahman is real. In the theistic schools, maya is not illusion but the real creative power, the \textbf{shakti}, of God.

\subsection{Two models of how the world relates to Brahman}

The schools split on whether the change from Brahman to world is only apparent or genuinely real.

\begin{longtable}{@{} L{0.18\textwidth}L{0.52\textwidth}L{0.22\textwidth}@{}}
\caption{The two models of the world's relation to Brahman.}\\
\toprule
\textbf{Model} & \textbf{Meaning} & \textbf{School} \\
\midrule
Vivarta & apparent transformation. Brahman appears as the world without truly changing, as a rope appears to be a snake in dim light. & Advaita (Shankara) \\
Parinama & real transformation. The cause genuinely becomes the effect, as milk becomes curd. The world is a real modification. & Samkhya, Vishishtadvaita, theistic schools \\
\bottomrule
\end{longtable}

\subsection{Ishvara and the Trimurti}

\textbf{Ishvara} is the \textbf{personal God}, the Lord, the absolute seen as creator, ruler, and object of devotion. In the theistic schools this is the true ultimate, identified with Vishnu or Shiva or the Goddess. The divine work is also expressed as three functions held by three figures (the \textbf{Trimurti}). \textbf{Brahma} the creator brings the worlds forth at the start of each cosmic day. \textbf{Vishnu} the preserver sustains them and descends as avatars to restore order when it slips. \textbf{Shiva} the destroyer and transformer dissolves the worlds at the close of the cycle, returning all to the unmanifest.

\subsection{The Advaita versus dualist split}

The schools disagree on which definition of the ultimate is primary. This is the central fault line of Hindu thought.

\begin{longtable}{@{} L{0.24\textwidth}L{0.22\textwidth}L{0.42\textwidth}@{}}
\caption{Where the schools place the ultimate.}\\
\toprule
\textbf{School} & \textbf{Stance} & \textbf{The ultimate} \\
\midrule
Advaita Vedanta (Shankara) & non-dualism & Nirguna Brahman alone is real. The world and the personal God are real only at the everyday level. \\
Vishishtadvaita (Ramanuja) & qualified non-dualism & Brahman is a personal God, Vishnu, with the world and the souls as his real body. \\
Dvaita (Madhva) & dualism & God and the souls are eternally distinct. \\
Samkhya & dualism & No God needed. Two eternal principles, consciousness and matter. \\
Bhakti, theistic & devotional & A personal supreme God is the highest reality. \\
\bottomrule
\end{longtable}

\subsection{The core teaching of the Bhagavad Gita}

The \textbf{Bhagavad Gita}, a dialogue set inside the great epic, is the most widely read single text of the tradition. On the surface it is a conversation on a battlefield, where Krishna counsels the warrior Arjuna. Read inwardly it is a treatise on the soul. Its core teachings draw the system together.

\textbf{The deathless self.} The first teaching is that the Atman never dies. "It is never born, and it never dies. Weapons do not cut it, fire does not burn it, water does not wet it, wind does not wither it." As a person sheds worn clothes for new ones, the self sheds worn bodies for new ones. The body is the garment. The self is the wearer.

\textbf{The field and the knower.} The Gita divides the embodied being in two. The \textbf{field}, called \textbf{kshetra}, is everything that can be observed, the body, the senses, the mind, the emotions, all of nature. The \textbf{knower of the field}, called \textbf{kshetrajna}, is the conscious witness for whom all of that is the seen. Bondage is the confusion of the two, the witness mistaking itself for what it merely watches. Liberation is the discrimination of the field from its knower.

\textbf{The three paths.} The Gita teaches several roads, fitted to different temperaments, all reaching the same goal.

\begin{longtable}{@{} L{0.26\textwidth}L{0.22\textwidth}L{0.40\textwidth}@{}}
\caption{The yogas, the paths of the Gita.}\\
\toprule
\textbf{Path} & \textbf{Native name} & \textbf{Means} \\
\midrule
path of action & Karma Yoga & selfless action with no attachment to its fruits \\
path of devotion & Bhakti Yoga & loving surrender to the personal God \\
path of knowledge & Jnana Yoga & discrimination and insight into the self \\
path of meditation & Dhyana Yoga & disciplined stilling of the mind \\
\bottomrule
\end{longtable}

The path of action is the Gita's most distinctive teaching. To be embodied is to act, so action cannot be stopped. The answer is not to drop action but to drop attachment to its results. Act fully and release the outcome, so that even forceful action leaves no residue on the soul.

\textbf{The vision of the universal form.} At the visionary center of the Gita, Krishna grants Arjuna a divine eye and reveals his \textbf{universal form}, a single boundless body holding all the worlds, with countless mouths and eyes and arms, the sun and moon as its eyes, its mouth a blazing fire. Arjuna sees all beings rushing into those mouths. He asks who this is. The answer is the most quoted line of the text. "I am Time, the destroyer of worlds." Birth and death are shown as two faces of one process. The form is described as opening only to undivided devotion.

\subsection{The split, restated}

Behind every one of these teachings runs the same divide. The non-dualist reads the deathless self as identical with the one Brahman, so the goal is \textbf{identity}, the drop discovering it was always the ocean. The dualist and qualified schools read the many selves as eternally real, each a distinct center of consciousness, so the goal is \textbf{loving communion}, the soul remaining itself yet joined to the Lord. The Gita holds both languages, of identity and of devotion, in one text.
